%!TEX program = xelatex

% 论文主文件：将 Markdown 转为 LaTeX 论文结构，保留原模板宏包配置
\documentclass[lang=cn,a4paper,11pt]{elegantpaper}
\usepackage{multirow,makecell,booktabs} % 表格美化、跨行
\usepackage{enumitem} % 列表格式
\usepackage{hyperref} % 超链接
\usepackage{geometry} % 页边距
\usepackage{fancyhdr}
\geometry{left=2.5cm,right=2.5cm,top=2.5cm,bottom=2.5cm}


\title{绿色财政与数字化治理政策对城市绿色发展的影响研究\\——基于两篇实证论文的整合分析}
\author{陶樞臣整理}
\institute{整理于 2026年}
\version{}
\date{2026年}

% 自定义列表：无缩进、紧凑
\setlist[enumerate]{leftmargin=2em, itemsep=0pt, parsep=0pt}
\setlist[itemize]{leftmargin=2em, itemsep=0pt, parsep=0pt}

\begin{document}
\maketitle

\begin{abstract}
本文整合两篇实证研究，分别考察绿色财政试点(节能减排财政政策综合示范城市)与政府数字化治理试点(信息惠民国家试点)对城市绿色发展的影响(见张静静等，2026；王青霞等，2026)。通过比较研究设计、变量测度、实证结果与机制分析，归纳两类政策在提升城市生态韧性、抑制能源强度与碳排放方面的共性与差异，并提出政策协同与实践建议。
\end{abstract}

\tableofcontents
\vspace{1em}

\section{引言}
本文整合两篇核心实证论文，分别梳理两项关键政策对城市绿色发展的影响，其一为“节能减排财政政策综合示范城市”试点，该政策属于绿色财政政策范畴；其二为“信息惠民国家试点”政策，其核心指向政府数字化治理。本文重点对比两篇研究的研究设计、实证结果与核心结论，为城市绿色发展政策实践提供综合参考。两篇论文均采用准自然实验与渐进双重差分模型开展实证研究，基于全国地级市面板数据展开分析，核心聚焦政策因果效应识别与多维度稳健性检验，二者在研究视角、核心变量设定与影响路径传导上存在显著差异。

\section{第一篇：绿色财政试点对城市生态韧性的影响}
\subsection{研究背景与研究目的}
城市生态韧性是支撑生态文明建设与城市可持续发展的核心基础，党的二十届三中全会明确提出推进城市生态韧性建设的核心要求。2011年，财政部与国家发改委联合推出节能减排财政政策综合示范城市试点，这项政策是我国绿色财政政策体系改革的核心试点与标杆实践，政策通过财政资金引导、刚性目标考核、多主体协同联动等方式，系统性推动城市产业低碳化转型、交通体系清洁化升级，搭建政府引导、市场主导、社会广泛参与的长效节能减排治理机制。本文选取2007至2023年我国285个地级市的面板数据作为研究样本，系统探究该项试点政策对城市生态韧性的影响效应、内在传导机制与差异化异质性特征，填补现有研究在多主体协同治理机制剖析、高维精准量化方法应用领域的研究空白。
\subsection{研究设计}
\subsubsection{基准模型设定}
为精准识别政策的因果效应，本文采用适配多批次、分时段试点的渐进双重差分模型作为基准回归模型，核心公式如下：
\[
Y_{it} = \alpha + \beta D_{it} + \gamma X_{it} + \delta T_{t} + \theta I_{i} + \epsilon_{it}
\]


公式中，下标i对应研究样本中的城市个体，下标t对应观测年份。被解释变量Resilience代表城市生态韧性水平，核心解释变量Policy为节能减排试点的虚拟变量，当城市i在t年份正式纳入示范城市试点范围时，该变量赋值为1，未纳入试点的年份与城市统一赋值为0。X代表全文选取的所有控制变量集合，α₀为模型截距项，α₁、α₂为模型待估核心系数，μᵢ为城市固定效应，λₜ为年份固定效应，εᵢₜ为模型随机扰动项。
\subsubsection{核心变量设定与测度}
第一类为被解释变量城市生态韧性。本文遵循学界通用的研究框架，从生态抵抗力、适应力、恢复力三个核心维度构建完整评价体系，指标覆盖工业污染排放压力、资源循环利用效率、生态修复保障能力、社会稳定支撑等多个细分领域，为提升指标测度的客观性与精准度，摒弃传统单一熵权法，采用熵权\-Critic组合赋权法完成综合指数测算，最大限度规避方差差异带来的权重偏误问题。

第二类为核心解释变量，即节能减排财政政策综合示范城市试点虚拟变量。本文依据国家部委正式发布的三批示范城市名单与官方批复时间完成变量赋值，研究样本共覆盖全国285个地级市与副省级城市，其中30个正式入选试点的城市构成处理组，剩余255个未入选试点的城市作为对照组，依托分批次试点的特征采用渐进双重差分模型完成因果识别。

第三类为机制传导变量，对应政策赋能生态韧性的三条核心路径。第一条路径为央地协同，本文采用空间向量夹角加权法，量化城市能源消费结构向低碳目标收敛的优化效应，以此衡量央地财政权责协同、资金统筹联动的治理效果；第二条路径为政企联动，通过耦合协调模型，分别测算政府端绿色治理刚性投入与企业端绿色投资、绿色创新行为的协同匹配程度，以此量化政企双向联动的治理效能；第三条路径为政府内驱，采用文本挖掘分析法，统计城市政府工作报告中环境保护相关关键词的频次占比，经过标准化处理后，衡量地方政府生态治理的内生动力与制度响应强度。

第四类为控制变量，本文结合现有主流文献的研究规范，选取6项核心控制变量，排除城市经济、产业、金融、开放等层面的干扰因素，分别为产业结构、人均金融规模、人均科技支出水平、对外开放程度、制造业集聚水平、城镇化率，所有变量均完成标准化与对数化处理，保证回归结果的稳健性。
\subsubsection{样本范围与数据来源}
本文研究样本为2007年至2023年我国285个地级市的平衡面板数据，总观测样本量为4845个。其中，城市生态韧性相关指标数据来源于《中国城市统计年鉴》、EPS数据库、各地级市官方统计公报与环境质量公报；试点虚拟变量依据国家发改委、财政部发布的节能减排财政政策综合示范城市正式名单整理得到；控制变量与机制变量的基础数据，来源于《中国城市统计年鉴》、EPS数据库、CSMAR数据库、CNRDS数据库及美国国家海洋和大气管理局公开数据。
\subsection{实证结果与检验分析}
\subsubsection{基准回归结果}
基准回归结果显示，无论模型是否加入控制变量，无论是否同时控制城市固定效应与年份固定效应，核心解释变量试点虚拟变量的回归系数均在百分之一的统计水平上显著为正，基准假设得到充分验证，节能减排财政政策试点能够显著提升城市生态韧性水平。
\subsubsection{全维度稳健性检验}
为保证核心结论的可靠性，本文开展多轮、全维度稳健性检验。首先通过事件研究法完成平行趋势检验，检验结果显示，政策正式实施之前，处理组与控制组的生态韧性变化趋势无显著差异，完全满足双重差分模型的核心适用前提。其次通过虚构处理组完成安慰剂检验，在试点城市范围内随机生成虚假实验组，重复500次回归模拟，结果显示虚假回归系数基本符合正态分布，绝大多数系数不具备统计显著性，且与真实回归系数距离较远，证明不可观测的随机因素不会对本文估计结果产生干扰。

同时，本文引入双重机器学习模型，通过套索回归、支持向量机、随机森林等多种高维算法，剥离模型中时变混淆变量与非线性关系的干扰，回归结果依旧验证核心结论稳健。针对试点城市选择可能存在的内生性问题，本文同步采用倾向得分匹配双重差分模型、工具变量法完成内生性处理，选取明朝驿站数量与时间趋势的交乘项作为工具变量，该工具变量既与当代城市行政执行能力、政策落地适配性高度相关，又不会直接影响当期城市生态韧性水平，满足外生性与相关性要求，回归结果依旧显著为正。

除此之外，本文采用广义随机森林模型完成多轮再估计，通过调整决策树迭代次数验证结果稳定性；采用Bacon分解法解构多时点双重差分模型的估计偏差，排除早批次试点城市作为晚批次试点城市对照组带来的估计污染问题；同时在模型中加入智慧城市试点、全面创新改革试点、气候适应型城市试点等同期政策虚拟变量，排除其他政策的干扰效应；剔除直辖市特殊样本、加入省份固定效应、更换被解释变量测度方法后，核心结论均未发生改变，充分证明本文研究结果具备极强的稳健性。

\subsubsection{传导机制检验}
本文采用温忠麟中介效应三步法完成机制检验，同步通过Sobel检验验证结果稳健，最终明确政策的三级传导路径。节能减排政策通过央地协同、政企联动、政府内驱三条路径，间接推动城市生态韧性提升，三条路径的效应存在明显差异。其中央地协同的传导效应最为显著，是政策赋能的核心渠道；政企联动的传导效应次之，是政策效能转化的关键载体；政府内驱的传导效应相对较弱，主要受地方财政收支约束、传统GDP导向政绩考核体系的双重限制，内生治理动力的释放空间不足。

\subsubsection{异质性分析结果}
本文从地理区位、资源禀赋、政策约束、城市规模、人口流动性多个维度展开异质性分析，结果显示政策赋能效果存在显著的空间与城市特征差异。从地理区位来看，政策正向效应在胡焕庸线以南城市高度显著，北部及沿线城市效应不显著；从资源禀赋与监管约束来看，政策对资源型城市的赋能效果优于非资源型城市，环境重点保护城市的政策效应显著优于非重点保护城市；从城市规模来看，政策对大规模城市与小规模城市均具备显著正向作用，效应差距极小；从人口流动性来看，政策效应在人口净流入城市高度显著，在人口净流出城市不具备统计显著性。

\subsection{研究结论与边际贡献}
本文核心研究结论为，节能减排财政政策综合示范城市试点，能够显著提升城市生态韧性水平，该因果效应经过多轮内生性处理与稳健性检验后依旧成立。政策通过央地协同、政企联动、政府内驱三级路径实现效能传导，三条路径的效应强度依次递减，同时政策赋能效果在不同区位、不同资源禀赋、不同人口流动特征的城市中存在显著异质性。

本文的边际学术贡献主要体现在三个方面。第一，优化了城市生态韧性的测度体系与测算方法，完善三维框架下的细分指标覆盖，采用熵权\-Critic组合赋权法提升测度精准度；第二，创新研究方法，引入双重机器学习、广义随机森林等高维算法，克服传统双重差分模型无法处理时变混淆变量、非线性关系的缺陷，实现政策效应的精准识别；第三，完善机制研究体系，构建央地协同、政企联动、政府内驱的完整传导逻辑链条，深度剖析绿色财政政策赋能生态韧性的内在机理，为相关领域研究提供了新的分析视角。


\section{第二篇：信息惠民国家试点政策的节能减排效应研究}
\subsection{研究背景与研究目的}
在“双碳”目标与“数字中国”建设背景下，政府数字化治理成为培育城市节能减排新动力的重要抓手。现有研究对政府数字化治理的节能减排效应存在争议，且多聚焦于技术设施供给端或生产端，未充分关注以公共服务优化与治理变革为核心的数字化模式的环境效应。“信息惠民国家试点”作为政府数字化治理的重要试点，依托政务大数据与平台整合，重塑政府治理与公众参与模式，其跨领域溢出的节能减排效应亟待检验。本文基于2006—2021年中国283个地级市面板数据，采用渐进DID模型，实证检验该政策的节能减排效应、调节效应与异质性特征。

\subsection{研究设计}
\subsubsection{基准模型}
采用渐进DID模型，核心聚焦政策对能源强度与碳排放强度的抑制效应，模型设定与第一篇逻辑一致，核心解释变量为信息惠民试点虚拟变量，被解释变量为能源强度、碳排放强度。

\subsubsection{核心变量}
(1)被解释变量：能源强度、碳排放强度，衡量城市节能减排水平。

(2)核心解释变量为信息惠民国家试点虚拟变量，记为Policy，该变量用于表征政府数字化治理带来的政策冲击。

(3)调节变量包含地方政府治理效率与公众关注度两项，用于检验这两个因素对政策节能减排效应的调节作用。

(4)控制变量方面，结合现有研究成果与本文研究主题，选取6类核心控制变量以确保回归结果的可靠性，具体包括：①经济发展水平，采用人均GDP对数衡量，用于反映城市经济基础对节能减排的支撑能力；②产业结构，以第二产业增加值占GDP比重为衡量标准，反映产业重型化程度对能源消耗与碳排放的影响；③人口规模，采用年末常住人口对数表示，用于控制人口集聚带来的能源需求差异；④环保投入，用环境保护支出占地方财政一般预算支出的比重衡量，体现地方政府对环保工作的重视程度；⑤科技水平，以人均科技支出对数为衡量指标，反映技术进步对节能减排的推动作用；⑥对外开放水平，采用进出口货物金额占GDP的比重表示，控制外资引入对产业环保水平的影响。

\subsubsection{数据来源与样本}

样本为2006—2021年中国283个地级市面板数据，剔除数据缺失严重的城市后，最终获得3782个观测值。数据来源具体如下：能源强度、碳排放强度数据来源于《中国能源统计年鉴》《中国环境统计年鉴》；核心解释变量根据国家发改委、财政部等部门发布的试点名单及批复时间整理；调节变量中，地方政府治理效率采用数据包络分析法测度，公众关注度通过百度指数“环境保护”相关关键词检索频次衡量；控制变量数据来源于《中国城市统计年鉴》、EPS数据库、各地级市统计公报及中国环境监测总站公开数据。

\subsubsection{调节变量}
\subsection{实证结果与稳健性检验}
\subsubsection{基准回归与稳健性检验}

信息惠民国家试点政策显著降低城市能源强度与碳排放强度，回归系数分别为\-0\.045、\-0\.214，且均在百分之一水平上显著，表明政策对城市节能减排具有显著促进作用。为验证结论可靠性，开展系列稳健性检验：平行趋势检验显示政策实施前各期虚拟变量系数均不显著，证明处理组与控制组满足平行趋势假设；安慰剂检验通过500次随机模拟，虚假系数均值接近0且不显著，排除非观测因素干扰；内生性处理采用“信息惠民试点城市数量的滞后一期”作为工具变量，通过两阶段最小二乘法缓解内生性，核心结论依然成立；替换被解释变量为单位GDP能耗、单位GDP碳排放，剔除直辖市、省会城市等特殊样本后，回归结果均与基准回归一致，进一步印证政策效应的稳健性。

\subsubsection{调节效应}
调节效应检验结果显示，地方政府治理效率每提升1单位，政策对能源强度与碳排放强度的抑制效应就分别增强0\.006和0\.031；公众关注度则仅对减排产生显著正向调节作用，对节能环节无显著影响。


\subsubsection{异质性分析}
政策的节能减排效能在非资源型城市与非老工业基地更为显著，体现出政策适配性特征。
\subsection{结论与边际贡献}
核心结论为，信息惠民国家试点政策作为政府数字化治理的重要实践，可显著促进城市节能减排，且其效应受地方政府治理效率与公众关注度的调节，异质性特征明显。


\section{两篇论文核心对比与总结}
\subsection{共性特征}
1\. 研究方法：均采用准自然实验设计与渐进DID模型，聚焦政策因果效应识别，通过系列稳健性检验保障结论可靠性；

2\. 研究样本：均以中国地级市面板数据为研究对象，覆盖较长时间跨度，具有较强的代表性与普适性；

3\. 政策导向：均聚焦城市绿色发展，探讨政策对生态环境与节能减排的正向影响，为政策优化提供实证参考。
\subsection{差异特征}
1\. 研究视角：第一篇聚焦绿色财政政策，核心是通过财政资金引导、目标考核与多主体协同，提升城市生态韧性的“抵抗\-适应\-恢复”三维能力，侧重生态系统的长期稳定与自我优化；第二篇聚焦政府数字化治理政策，核心是通过政务大数据整合、公共服务优化与公众参与渠道拓宽，直接抑制能源强度与碳排放强度，侧重节能减排的短期量化成效，两者形成“长期生态赋能”与“短期减排增效”的互补视角。

2\. 核心变量：被解释变量方面，第一篇采用多维度综合指标，即城市生态韧性，该指标涵盖工业污染、资源循环、生态修复等18项细分指标，采用熵权\-Critic组合赋权法完成测度；第二篇采用单一维度核心指标，具体为能源强度与碳排放强度，可直接量化节能减排成效。机制与调节变量方面，第一篇侧重协同机制的挖掘，重点分析政策通过央地、政企、政府内驱三条路径的传导逻辑；第二篇侧重调节效应的探究，重点分析地方政府治理效率与公众关注度对政策效能的放大或抑制作用。

3\. 研究重点与方法细节：第一篇侧重政策对生态系统整体韧性的长期赋能，除采用渐进双重差分模型外，还引入双重机器学习、广义随机森林等高维算法，有效解决时变混淆变量与非线性关系带来的估计偏差问题；第二篇侧重政策对节能减排核心指标的短期抑制效应，重点分析调节机制，采用数据包络分析法测度地方政府治理效率，结合百度指数量化公众关注度，研究方法更侧重宏观调节效应的识别与检验。

\subsection{综合启示}

绿色财政政策与政府数字化治理政策均是推动城市绿色发展的重要抓手，两者可协同发力、互补增效：绿色财政政策通过资金引导与多主体协同，夯实城市生态韧性基础，为节能减排提供长期生态支撑；数字化治理政策通过精准监测、高效协同与公众参与，提升节能减排效能，为生态韧性建设提供技术与治理保障。未来政策实践中，应注重两类政策的适配与融合：一是推动绿色财政资金向数字化环保领域倾斜，支持城市生态监测数字化平台建设，提升财政资金使用效率；二是利用数字化手段优化绿色财政政策执行，通过大数据精准识别政策适配城市与重点领域，避免资金浪费；三是结合城市异质性特征，针对胡焕庸线以南城市、环境重点保护城市，强化两类政策的协同力度，针对资源型城市、人口净流出城市，优化政策侧重点，充分释放政策协同效应；四是完善配套机制，将节能减排成效与生态韧性提升纳入地方政府政绩考核，强化政府内驱动力，同时拓宽公众参与渠道，形成“政策引导、技术支撑、公众参与”的绿色发展格局，助力“双碳”目标实现与城市可持续发展。

\section*{参考文献}

\begin{thebibliography}{9}
\bibitem{Zhang2026} 张静静, 谢佳辉, 刘逸民, 等. 绿色财政政策如何赋能城市生态韧性？——来自“节能减排示范城市”准自然试验的证据. 中央财经大学学报, 2026(3):17--34.
\bibitem{Wang2026} 王青霞, 王巧, 李斌. 信息惠民国家试点政策的节能减排效应研究. 环境科学研究, 2026,39(2):1--12.
\bibitem{Wang2025} 王伟成. 城市数字治理与新质生产力——来自“信息惠民国家试点”政策的证据. 建筑与文化, 2025(6):266--269.
\bibitem{He2024} 何雨可, 牛耕, 逯建, 等. 数字治理与城市创业活力——来自“信息惠民国家试点”政策的证据. 数量经济技术经济研究, 2024,41(1):47--66.
\bibitem{MOF2011} 财政部, 国家发展和改革委员会. 关于开展节能减排财政政策综合示范工作的通知. 2011.
\bibitem{NDRC2014} 国家发展和改革委员会, 财政部, 工业和信息化部, 等. 关于开展信息惠民国家试点工作的通知. 2014.
\bibitem{StatYearbook} 中国统计出版社. 中国城市统计年鉴(2007--2023). 北京: 中国统计出版社, 2007--2023.
\bibitem{EnvYearbook} 中国环境科学出版社. 中国环境统计年鉴(2006--2021). 北京: 中国环境科学出版社, 2006--2021.
\bibitem{Wen2004} 温忠麟, 张雷, 侯杰泰, 等. 中介效应检验程序及其应用. 心理学报, 2004,36(5):614--620.
\end{thebibliography}

\end{document}